\Article{9}
Just as the preceding method is based on dividing the region $S$ into prismatic elements, the second method will proceed by dividing the same region into pyramidal elements. Imagine a spherical surface of radius $=1$ drawn about the center $\mu$ and divided into infinitely small elements. Toward a point of such an element $\mathrm{d}\Pi$ draw a straight line from the point $\mu$. Let it cut the surface of the region $S$ successively at the points $P',P'',P'''$ etc. Let the distances of these points from $\mu$ be denoted by $R',R'',R'''$ etc. The straight lines drawn from $\mu$ toward all points of the boundary of the element $\mathrm{d}\Pi$ will form a pyramidal region and will cut out surface elements at the points $P',P'',P'''$ etc. of the region $S$, which we denote by $\mathrm{d}T',\mathrm{d}T'',\mathrm{d}T'''$ etc. Finally, let $Q'$ be the angle between the straight line $P'\mu$ and the outward normal to the element $\mathrm{d}T'$. Likewise, let $Q'',Q'''$ etc. be the inclinations of the corresponding normals at the points $P'',P'''$ etc. to the straight line drawn toward $\mu$. Then
\[
\mathrm{d}\Pi=\pm\frac{\mathrm{d}T'.\cos Q'}{R'R'}=\mp\frac{\mathrm{d}T''.\cos Q''}{R''R''}=\pm\frac{\mathrm{d}T'''.\cos Q'''}{R'''R'''}\ \text{etc.}
\]
where the upper or lower signs apply according as $\mu$ lies outside or inside the region $S$. The case in which $\mu$ lies on the surface of the region $S$ itself is to be counted with the first or the second case according as the line $\mu P'$ falls outside or inside the region $S$.

Further, it is clear that for all parts of the region $S$ within that pyramidal region, the angle $q$ is constant. Thus, by reasoning like that in Article 7, we conclude that, if we set the indefinite integral
\[
\int\psi r.\mathrm{d}r=-\theta r
\]
with the constant of integration chosen at will, the integral
\[
\int\frac{\mathrm{d}S.\cos q.\psi(\mu,\mathrm{d}S)}{(\mu,\mathrm{d}S)^2}
\]
taken over all parts of the region $S$ within that pyramidal region, will be in the first case
\SourceLeaf{48}
\[
=\cos q.\left(\frac{\mathrm{d}T'.\cos Q'.\theta R'}{R'R'}+\frac{\mathrm{d}T''.\cos Q''.\theta R''}{R''R''}+\frac{\mathrm{d}T'''.\cos Q'''.\theta R'''}{R'''R'''}+\ \text{etc.}\right)
\]
while in the second case the term $\mathrm{d}\Pi.\cos q.\theta0$ must be added to the same formula.

Now, on summing over all elements of the spherical surface, the complete integral
\[
\int\frac{\mathrm{d}S.\cos q.\psi(\mu,\mathrm{d}S)}{(\mu,\mathrm{d}S)^2}
\]
will be

I. in the case in which the point $\mu$ lies outside the region $S$,
\[
=\int\frac{\mathrm{d}T.\cos q.\cos Q.\theta R}{RR}
\]
where $\mathrm{d}T$ denotes any element of the surface of the region $S$, and $Q,R$ denote for these elements the same quantities as the primed letters did earlier for particular elements. Finally, $q$ is the inclination to our fixed line of the straight line drawn from the point $\mu$ toward the element $\mathrm{d}T$.

II. In the case in which the point $\mu$ lies inside the region $S$, the term
\[
\theta0.\int\mathrm{d}\Pi.\cos q
\]
must be added, where $q$ is the inclination of the straight line drawn from $\mu$ toward $\mathrm{d}\Pi$ to the fixed line, and the integration must extend over the entire spherical surface. It is easy to see, however, that this integral taken over the hemisphere on which $q$ is acute is $=+\pi$, and over the other hemisphere it is $=-\pi$. Hence the complete integral vanishes, and exactly the same formula we gave for the first case holds for this second case. The situation is different in the third case.

III. When the point $\mu$ lies on the surface of the region $S$ itself. Here too the term
\[
\theta0.\int\mathrm{d}\Pi.\cos q
\]
must be added, but the integration is to extend only over those parts of the spherical surface for which the initial portion of the straight line drawn from $\mu$ toward $\mathrm{d}\Pi$ lies within the region $S$. Equivalently, provided the surface of the region $S$ at the point $\mu$ has neither a cusp nor an edge, these are the parts for which this line makes an obtuse angle with the outward normal to the surface of the region $S$ at the point $\mu$. It remains, then, to find the integral understood in this sense.

\SourceLeaf{49}Let this normal and the fixed line cut the spherical surface at the points $G,H$, respectively. Let the arc $GH=k$, and let the arc between $G$ and an arbitrary point of the spherical surface be $=v$. Finally, let $w$ be the spherical angle between the arcs $k,v$. Then $\cos q=\cos k.\cos v+\sin k.\sin v.\cos w$, and for $\mathrm{d}\Pi$ we must take the element $\sin v.\mathrm{d}v.\mathrm{d}w$. The integral $\int\mathrm{d}\Pi.\cos q$
\[
=\int\int(\cos k.\cos v+\sin k.\sin v.\cos w)\sin v.\mathrm{d}v.\mathrm{d}w
\]
must extend from $w=0$ to $w=360^0$ and from $v=90^0$ to $v=180^0$. In this way the first integration gives
\[
\int2\pi\cos k.\cos v.\sin v.\mathrm{d}v
\]
and then the second gives $-\pi\cos k$.

For our purpose, this third case comes into consideration only insofar as the surfaces of the regions $s,S$ have a finite part in common. If the point $\mu$ lies on that part, then either $k=0$ or $=180^0$, so the integral $\int\mathrm{d}\Pi.\cos q$ is either $=-\pi$ or $=+\pi$, according as, at the point $\mu$, the regions $s,S$ lie on the same side or on opposite sides of the plane tangent to both surfaces.

Applying this to our original integral $\int\int\mathrm{d}s.\mathrm{d}S.\varphi(\mathrm{d}s,\mathrm{d}S)$, its value becomes

I. whenever the surfaces of the regions $s,S$ have no finite part in common,
\[
=4\pi\sigma\psi0+\int\int\frac{\mathrm{d}t.\mathrm{d}T.\cos q.\cos Q.\theta(\mathrm{d}t,\mathrm{d}T)}{(\mathrm{d}t,\mathrm{d}T)^2}
\]

II. whenever the surfaces of the regions $s,S$ have a finite part $=\tau$ in common,
\[
=4\pi\sigma\psi0\mp\pi\tau\theta0+\int\int\frac{\mathrm{d}t.\mathrm{d}T.\cos q.\cos Q.\theta(\mathrm{d}t,\mathrm{d}T)}{(\mathrm{d}t,\mathrm{d}T)^2}
\]
where the upper or lower sign applies according as the regions $s,S$ lie on the same side or on opposite sides of their common surface $\tau$.

III. Whenever the surfaces of the regions $s,S$ have several separate finite parts in common, let $\tau$ be the sum of those that the regions $s,S$ adjoin from the same side, and $\tau'$ the sum of those that these regions adjoin from opposite sides. Our integral will then be
\SourceLeaf{50}
\[
=4\pi\sigma\psi0+\pi(\tau'-\tau)\theta0+\int\int\frac{\mathrm{d}t.\mathrm{d}T.\cos q.\cos Q.\theta(\mathrm{d}t,\mathrm{d}T)}{(\mathrm{d}t,\mathrm{d}T)^2}
\]
This third formula can be regarded as covering all cases. The double integral must extend over all elements of both surfaces, and $q,Q$ denote the angles that the straight line joining the two elements $\mathrm{d}t,\mathrm{d}T$ makes with the outward normals to these elements, taking the direction of the line from $\mathrm{d}t$ toward $\mathrm{d}T$ for the first angle and from $\mathrm{d}T$ toward $\mathrm{d}t$ for the second.

\Article{10}
The two transformations of the integral $\int\mathrm{d}s.\mathrm{d}S.\varphi(\mathrm{d}s,\mathrm{d}S)$ developed in Articles 8 and 9 recommend themselves by their nearly equal elegance, though the latter is better suited to our purpose. The general problem cannot be reduced further unless we make specific assumptions about either the regions $s,S$ or the function $\varphi$. And since the function $\varphi$ arises from the function $f$, we shall base the rest of our investigation on the same hypothesis from which the distinguished \textsc{Laplace} started, namely that the molecular forces of attraction are perceptible only at imperceptible distances. Since this expression is somewhat vague as long as no unit is specified, we note first of all that the attractive force $fr$, expressed as a function of distance $r$, must first be multiplied by a mass to have the same dimensions as gravity $g$. The meaning of the assumption is then that, with $M$ denoting a mass such as might occur in experiments, one that may be treated as negligible in comparison with the entire Earth, $Mfr$ must always remain imperceptible in comparison with gravity as long as $r$ has a value our measurements can register, however small. Nothing prevents the value of $Mfr$ from becoming perceptible at imperceptible distances or even from exceeding all bounds as $r$ decreases. We cannot help admiring the importance and quite distinctive mathematical character of the results that can be drawn from this one hypothesis, while the law of the function $fr$ is otherwise regarded as entirely unknown. Under these circumstances they can lay no claim to absolute mathematical precision, yet they certainly retain such precision that no experiment can detect any departure from absolute truth. For as soon as we succeeded in subjecting such a departure to measurement, the assumption itself would cease to hold.

\SourceLeaf{51}\Article{11}
We may assume that the function $fr$, and likewise the function $Fr$, denotes the attraction with the part inversely proportional to $rr$ that serves to explain astronomical phenomena left out. Whatever the shapes of the fluid and the vessel, this part can cause only an imperceptible change in gravity at any point. Thus, as $r$ increases from a perceptible value to infinity, $fr$ will itself be imperceptible and will also decrease faster than $\frac{1}{rr}$. It is easy to infer from this that the integral $\int fr.\mathrm{d}r$ taken from any perceptible value to infinity is also imperceptible. We shall therefore assume that the constant of integration in $\int fr.\mathrm{d}r=-\varphi r$ is chosen so that $\varphi\infty=0$, or so that $\varphi r$ is the value itself of the integral $\int fx.\mathrm{d}x$ taken from $x=r$ to $x=\infty$. In this way $\varphi r$ denotes a positive quantity for any distance $r$, though imperceptible as long as $r$ is perceptible. For an imperceptible value of $r$, on the other hand, it may become perceptible and even exceed all bounds as the distance $r$ continues to decrease. In ordinary language, nothing prevents $\varphi0=\infty$.

